\documentclass[../cours_1re.tex]{subfiles}
\begin{document}
\setcounter{chapter}{28}
\chapter{Complément : convergence des algorithmes}
\label{ch:29}
\begin{objectif}
Complément : relier une itération à sa convergence, à sa vitesse et à ses erreurs.
Prérequis : suites, dérivation, exponentielle et complétude du chapitre 25.
\end{objectif}
\section{Une contraction possède un point fixe}
Soit $I$ un intervalle fermé non vide, $f(I)\subset I$ et $0\le q<1$ tels que
$|f(x)-f(y)|\le q|x-y|$ pour tous $x,y\in I$.
La suite $u_{n+1}=f(u_n)$, $u_0\in I$, reste dans $I$.
Par récurrence, $|u_{n+1}-u_n|\le q^n|u_1-u_0|$.
Pour $m>n$, l'inégalité triangulaire et la somme géométrique donnent
\[
|u_m-u_n|\le\frac{q^n}{1-q}|u_1-u_0|.
\]
La suite est de Cauchy, donc converge vers $\ell\in I$.
La condition lipschitzienne implique la continuité, donc $f(\ell)=\ell$.
Deux points fixes $\ell,\ell'$ vérifieraient
$|\ell-\ell'|\le q|\ell-\ell'|$ ; puisque $q<1$, ils sont égaux.
En passant à la limite dans la borne précédente :
\[
|u_n-\ell|\le\frac{q^n}{1-q}|u_1-u_0|.
\]
Après avoir calculé $u_n$ pour $n\ge1$, on dispose aussi de
$|u_n-\ell|\le q|u_n-u_{n-1}|/(1-q)$.
Pour $q=0$, la fonction est constante et la solution est atteinte dès le
premier pas ; les formules s'interprètent directement, sans ambiguïté de $0^0$.

Si $f$ est dérivable sur $I$ et $|f'|\le q$, le théorème des accroissements
finis prouve la condition de contraction. Vérifier aussi $f(I)\subset I$ :
une dérivée petite ne suffit pas si les itérés quittent le domaine.

\section{Newton : une vitesse quadratique, sous conditions}
Pour la racine $r=\sqrt A$, $A>0$, Newton donne
$u_{n+1}=(u_n+A/u_n)/2$. Si $u_0>0$, $u_1\ge r$ puisque
$u_1-r=(u_0-r)^2/(2u_0)\ge0$.
Ensuite $u_n\ge r$, $u_{n+1}\le u_n$ et la limite est $r$
(chapitre 11). L'erreur vérifie exactement
\[
0\le u_{n+1}-r=\frac{(u_n-r)^2}{2u_n}
\le\frac{(u_n-r)^2}{2r}.
\]
L'erreur est donc essentiellement mise au carré à chaque étape près de la solution.
L'intervalle $[A/u_n,u_n]$ contient $r$ et fournit un arrêt justifié.
Cette garantie porte sur les valeurs mathématiques exactes. Avec des flottants,
il faut tenir compte des arrondis pour obtenir un encadrement numériquement certifié.
Pour une fonction générale, des hypothèses locales supplémentaires sont
nécessaires ; l'exemple cyclique du chapitre 11 interdit d'annoncer
une convergence universelle de Newton.

\section{Euler : de l'erreur locale à l'erreur globale}
Approchons $y'=y$, $y(0)=1$, sur $[0,T]$, $T>0$, avec $h=T/n$.
Les valeurs d'Euler sont $Y_k=(1+h)^k$.
La formule de Taylor avec reste du chapitre 11 donne
\[
e^{(k+1)h}=e^{kh}+he^{kh}+r_k,\qquad
0\le r_k\le\tfrac12e^Th^2.
\]
En posant $d_k=e^{kh}-Y_k$, on a
$d_{k+1}=(1+h)d_k+r_k$ et $d_0=0$.
En déroulant la récurrence :
\[
0\le d_n\le\tfrac12e^Th^2\sum_{j=0}^{n-1}(1+h)^j
=\tfrac12e^Th\big((1+h)^n-1\big)
\le\tfrac12e^T(e^T-1)h.
\]
La dernière inégalité vient de $1+h\le e^h$.
L'erreur globale est donc au plus une constante fois $h$ :
Euler converge avec un ordre au moins un dans cet exemple.
L'erreur locale est en $h^2$, mais il y a environ $T/h$ étapes.
On retrouve $(1+1/n)^n\to e$ pour $T=1$.

\section{Stabilité et choix du pas}
Pour $y'=-ky$, $k>0$, Euler donne $Y_{j+1}=(1-kh)Y_j$.
Il tend vers $0$ si $|1-kh|<1$, c'est-à-dire $0<kh<2$.
Pour $0<kh\le1$, les valeurs positives restent positives.
Pour $1<kh<2$, elles alternent de signe tout en décroissant en valeur absolue,
contrairement à la solution exacte positive. Pour $kh>2$, elles divergent
en valeur absolue. Convergence d'une méthode lorsque $h\to0$ et stabilité
pour un pas fixé sont deux questions distinctes.

La descente de gradient pour minimiser
$Q(x)=\frac a2(x-b)^2$, $a>0$, donne
$x_{n+1}=x_n-\alpha Q'(x_n)$, donc
$x_{n+1}-b=(1-\alpha a)(x_n-b)$.
La même analyse donne convergence pour $0<\alpha a<2$.
Une méthode d'optimisation se comprend ici par une suite géométrique.

\section{Conditionnement et calcul flottant}
Résoudre $x^2=\varepsilon$ donne $x=\sqrt\varepsilon$ :
près de $0$, une petite variation absolue du second membre peut être
fortement amplifiée relativement. C'est une sensibilité du problème.
À cela peut s'ajouter une instabilité du calcul : dans
$(-b+\sqrt{b^2-4ac})/(2a)$, soustraire deux grands nombres proches perd
des chiffres significatifs.
Lorsque $b\ne0$ et $\Delta>0$, calculer d'abord
$q=-\frac12(b+\operatorname{sgn}(b)\sqrt\Delta)$, puis les racines
$q/a$ et $c/q$. La seconde formule vient du produit des racines $c/a$.
Traiter séparément les cas où ces divisions seraient impossibles.
Cette réécriture réduit une source d'annulation ; elle ne corrige pas
la sensibilité intrinsèque de racines presque confondues.

\begin{finchapitre}
\section{Utilité et prolongements}
Une formule d'erreur transforme un critère d'arrêt en garantie.
Les mêmes suites géométriques contrôlent contractions, stabilité d'Euler
et descente de gradient sur une fonction quadratique.

L'analyse numérique et l'optimisation de master distingueront systématiquement
existence, convergence, vitesse, conditionnement et coût. Ces questions
doivent être posées avant de faire confiance à un résultat informatique.
\begin{retenir}
Une itération ne converge pas par définition. Prouver la stabilité du domaine,
établir une estimation d'erreur et distinguer les limites du modèle de celles du calcul.
\end{retenir}
\end{finchapitre}
\end{document}
